% ============================================================
%  TEMA 3 — APARTADO 2: ANÁLISIS DE OFERTA Y DEMANDA DE RRHH
%
%  RECORDATORIO AL EQUIPO:
%
%  ANÁLISIS DE DEMANDA (¿cuántos empleados necesitamos?):
%  Métodos cuantitativos:
%    - Análisis de series temporales (extrapolación de tendencias históricas).
%    - Ratios de productividad (unidades producidas / empleado).
%    - Análisis de regresión (demanda de empleo en función de variables del negocio).
%    - Modelos econométricos.
%  Métodos cualitativos:
%    - Técnica Delphi (consenso de expertos internos).
%    - Juicio directivo (managers estiman sus necesidades).
%    - Escenarios (qué pasa si... → planificación por escenarios).
%
%  ANÁLISIS DE OFERTA (¿cuántos empleados tenemos / tendremos?):
%  Oferta interna (inventario de la plantilla actual):
%    - Inventario de habilidades / perfiles.
%    - Matrices de Markov / tablas de transición (probabilidad de que un
%      empleado permanezca, sea promovido, se vaya...).
%    - Análisis de la rotación y del absentismo.
%    - Mapas de reemplazamiento / planes de sucesión.
%  Oferta externa (mercado laboral):
%    - Análisis del mercado de trabajo local / sectorial.
%    - Datos del INE, SEPE, convenios colectivos.
%    - Benchmarking con otras entidades financieras.
%
%  PARA ESTA SECCIÓN INDICAD:
%  - ¿Qué métodos de análisis de demanda usa CaixaBank?
%  - ¿Qué métodos de análisis de oferta usa?
%  - ¿Cuáles serían convenientes aplicar que no usa actualmente?
%  - Ventajas e inconvenientes del enfoque actual.
% ============================================================

\section{Análisis de oferta y demanda de recursos humanos}

\subsection{Análisis de la demanda de RRHH}

La previsión de la demanda de personal permite a CaixaBank anticipar cuántos
y qué tipo de empleados necesitará en el futuro para alcanzar sus objetivos
de negocio.

\subsubsection{Métodos utilizados para la previsión de la demanda}

\begin{table}[H]
  \centering
  \small
  \caption{Métodos de previsión de la demanda de RRHH en CaixaBank}
  \begin{tabularx}{\textwidth}{p{0.28\textwidth} p{0.12\textwidth} X}
    \toprule
    \textbf{Método}                  & \textbf{¿Se usa?} & \textbf{Descripción y observaciones} \\
    \midrule
    Series temporales / tendencias   & [Sí/No/Parcial] & \\
    Ratios de productividad          & [Sí/No/Parcial] & \\
    Regresión / modelos cuantitativos & [Sí/No/Parcial] & \\
    Técnica Delphi                   & [Sí/No/Parcial] & \\
    Juicio directivo                 & [Sí/No/Parcial] & \\
    Planificación por escenarios     & [Sí/No/Parcial] & \\
    \bottomrule
  \end{tabularx}
\end{table}

[Descripción detallada de los métodos más relevantes que usa CaixaBank]

\textbf{Métodos que serían convenientes aplicar:}

% TODO: Basándoos en el tamaño y complejidad de CaixaBank, ¿qué métodos
% más avanzados podrían añadir valor? Ejemplos:
%   - Modelos de regresión múltiple incorporando variables macroeconómicas
%     (volumen de crédito, tipos de interés, PIB) para prever necesidades.
%   - People Analytics con machine learning para predecir rotación.
%   - Planificación por escenarios ante disrupciones tecnológicas (IA, blockchain).

\begin{itemize}
  \item [Método 1 — justificación]
  \item [Método 2 — justificación]
\end{itemize}

\subsection{Análisis de la oferta de RRHH}

\subsubsection{Oferta interna: inventario de la plantilla}

% TODO: ¿Tiene CaixaBank un sistema de inventario de habilidades/competencias?
% ¿Usan matrices de Markov o tablas de transición para modelar flujos internos?
% ¿Tienen mapas de reemplazamiento / planes de sucesión formales?

[Descripción a completar]

\subsubsection{Oferta externa: análisis del mercado de trabajo}

% TODO: ¿Cómo monitoriza CaixaBank el mercado laboral financiero en España?
% ¿Hacen benchmarking de salarios con competidores (Santander, BBVA, Sabadell)?
% ¿Colaboran con universidades para captar talento (relaciones con la UGR,
% prácticas, campus recruitment)?

[Descripción a completar]

\begin{notaequipo}
  Preguntad en la entrevista:
  \begin{enumerate}
    \item ¿Con qué métodos cuantifican sus necesidades futuras de personal?
    \item ¿Tienen un inventario actualizado de competencias de la plantilla?
    \item ¿Usan modelos matemáticos o estadísticos para prever la oferta interna
          (p.ej. matrices de transición / Markov)?
    \item ¿Cómo monitorizan el mercado laboral externo?
    \item ¿Cuál es su tasa de rotación actual? ¿La consideran alta o baja?
    \item ¿Han tenido dificultades para cubrir ciertos perfiles en los últimos años?
          ¿Cuáles? (Probablemente perfiles tecnológicos: Data Scientists, Ingenieros,
          especialistas en ciberseguridad)
  \end{enumerate}
\end{notaequipo}
